% !TEX root = 第三章_一元函数的连续性(答案版).tex
\setcounter{chapter}{2}
\pagestyle{fancy}

\chapter{一元函数的连续性}


\section{连续性证明}

要证明一个函数 $f$ 在某区间 $I$ 上连续,只要在区间里任意取定一点 $x_0\in I$,证明 $\displaystyle\lim\limits_{x\to x_0}f(x)=f(x_0)$.结合前面的极限知识,有以下几种方法证明:
\begin{enumerate}
    \item 利用定义---对任意 $\varepsilon>0$,存在 $\delta>0$,当 $|x-x_0|<\delta$ 时,有 $|f(x)-f(x_0)|<\varepsilon$;
    \item 利用左,右极限---$f(x_0+0)=f(x_0)=f(x_0-0)$;
    \item 利用数列极限(Heine 归结原理)---对任意 $\{x_n\}\to x_0$,有 $f(x_n)\to f(x_0)$;
    \item 利用连续函数的运算性质---连续函数的四则运算,复合下仍连续(即如下的两个定理).
\end{enumerate}

\begin{theorem}[\markstar]
    设 $f(x),g(x)$ 均在 $x_0$ 点连续,则:
    \begin{enumerate}
        \item $f(x)\pm g(x)$ 以及 $f(x),g(x)$ 均在点 $x_0$ 处连续;
        \item 当 $g(x_0)\neq0$ 时,$\dfrac{f(x)}{g(x)}$ 在点 $x_0$ 处连续.
    \end{enumerate}
\end{theorem}

\begin{theorem}[\markstar]
    设 $u=g(x)$ 在 $x_0$ 处连续,$y=f(u)$ 在点 $u_0=g(x_0)$ 处连续,则复合函数 $f(g(x))$ 在点 $x_0$ 处连续.
\end{theorem}

%例题3.1.1
\begin{example}
    研究函数
    \[
        f(x)=\lim\limits_{n\to\infty}\dfrac{x^n-1}{x^n+1}
    \]
    的连续性.
\end{example}
\exampleblank

\begin{solution}
    当 $|x|<1$ 时,$x^n\to0$;当 $|x|>1$ 时,商的分子分母同除以 $x^n$ 可得极限为 $1$.又 $x=1$ 时极限为 $0$,$x=-1$ 时原式对奇数 $n$ 无定义.故极限函数在其定义域上为
    \[
        f(x)=
        \begin{cases}
            1,  & x<-1\text{ 或 }x>1, \\
            -1, & -1<x<1,            \\
            0,  & x=1.
        \end{cases}
    \]
    因而它在 $x=1$ 间断,在其余定义点连续;$x=-1$ 不在定义域内.
\end{solution}

%例题3.1.2
\begin{example}[\markstar]
    证明 Riemann 函数
     \[
        R(x)=
        \begin{cases}
            \dfrac{1}{q}, & x=\dfrac{p}{q}\in(0,1),\ p,q\in\mathbb{N},\ (p,q)=1, \\
            0,            & x\in\left(\lbrack0,1\rbrack\setminus\mathbb{Q}\right)\cup\{0,1\},
        \end{cases}
    \]
    在无理点上连续,在有理点上间断.
\end{example}
\exampleblank

\begin{solution}
    设 $x_0$ 为无理数.给定 $\varepsilon>0$,只有有限多个分母 $q<\dfrac{1}{\varepsilon}$ 的既约分数落在任意有界区间中.取 $x_0$ 的一个充分小邻域避开这些点,则其中有理点 $\dfrac{p}{q}$ 均满足 $\dfrac{1}{q}<\varepsilon$,无理点上函数值为 $0$,故 $R(x)\to0=R(x_0)$.

    若 $x_0=\dfrac{p_0}{q_0}$ 为既约分数,取无理数列 $x_n\to x_0$,则 $R(x_n)=0$,但 $R(x_0)=\dfrac{1}{q_0}>0$,故在 $x_0$ 间断.
\end{solution}

%例题3.1.3
\begin{example}
    讨论函数
    \[
        f(x)= \begin{cases} x(1-x), & x\text{ 为有理数}, \\ x(1+x), & x\text{ 为无理数}, \end{cases}
    \]
    的连续性.
\end{example}
\exampleblank

\begin{solution}
    分别沿有理数列和无理数列趋于 $x_0$,所得极限分别为 $x_0(1-x_0),\,x_0(1+x_0)$.
    二者相等当且仅当 $x_0=0$.此时 $|f(x)|\leq |x|(1+|x|)\to0=f(0)$,故 $f$ 仅在 $x=0$ 连续.
\end{solution}

%例题3.1.4
\begin{example}
    设
    \[
        f(x)= \begin{cases} 1,  & x\geq0, \\ -1, & x<0, \end{cases} \qquad g(x)=\sin x,
    \]
    讨论 $f(g(x))$ 的连续性.
\end{example}
\exampleblank

\begin{solution}
    $f(g(x))$ 在 $\sin x\geq0$ 时等于 $1$,在 $\sin x<0$ 时等于 $-1$.在每个 $x=k\pi$ 的两侧,$\sin x$ 变号,因而左右极限不相等.故它在 $k\pi$ 处间断,在其余各点连续.
\end{solution}

%例题3.1.5
\begin{example}
    设 $f(x)$ 在 $(0,1)$ 内有定义,且函数 $e^x f(x)$ 与 $e^{-f(x)}$ 在 $(0,1)$ 内都是单调不减的,试证:$f(x)$ 在 $(0,1)$ 内连续.
\end{example}
\exampleblank

\begin{solution}
    由 $e^{-f(x)}$ 单调不减知 $f$ 单调不增,故任一点 $x_0$ 的左右极限存在,且 $f(x_0-)\geq f(x_0)\geq f(x_0+)$.
    另一方面,$e^xf(x)$ 单调不减,令 $x\to x_0-$ 和 $x\to x_0+$ 得
    \[
        e^{x_0}f(x_0-)\leq e^{x_0}f(x_0)\leq e^{x_0}f(x_0+).
    \]
    两组不等式迫使三个值相等,故 $f$ 在每点连续.
\end{solution}

%例题3.1.6
\begin{example}
    设 $f(x)$ 在 $(a,b)$ 上至多只有第一类间断点,且对任意 $x,y\in(a,b)$,
    \[
        f\left(\dfrac{x+y}{2}\right)\leq\dfrac{f(x)+f(y)}{2}.
    \]
    证明:$f(x)$ 在 $(a,b)$ 上连续.
\end{example}
\exampleblank

\begin{solution}
    任取 $x_0\in(a,b)$.由于 $x_0$ 至多为第一类间断点,故有限的单侧极限 $f(x_0-0),f(x_0+0)$ 均存在.对充分小的 $h\neq0$,使 $x_0\pm h\in(a,b)$,由题设分别有
    \[
        f(x_0)\leq\dfrac{f(x_0+h)+f(x_0-h)}{2},
        \qquad
        f\left(x_0+\dfrac{h}{2}\right)\leq\dfrac{f(x_0)+f(x_0+h)}{2}.
    \]
    在第一个不等式中令 $h\to0^+$,在第二个不等式中分别令 $h\to0^+$ 与 $h\to0^-$,得到
    \[
        \begin{aligned}
            f(x_0)&\leq\dfrac{f(x_0+0)+f(x_0-0)}{2},\\
            f(x_0+0)&\leq\dfrac{f(x_0)+f(x_0+0)}{2},\\
            f(x_0-0)&\leq\dfrac{f(x_0)+f(x_0-0)}{2}.
        \end{aligned}
    \]
    后两式分别给出 $f(x_0+0)\leq f(x_0)$ 与 $f(x_0-0)\leq f(x_0)$.代入第一式可得
    \[
        2f(x_0)\leq f(x_0+0)+f(x_0-0)\leq2f(x_0),
    \]
    因而
    \[
        f(x_0-0)=f(x_0)=f(x_0+0).
    \]
    故 $f$ 在 $x_0$ 连续.由 $x_0$ 的任意性,$f$ 在 $(a,b)$ 上连续.
\end{solution}

%例题3.1.7
\begin{example}[\markstar]
    \begin{enumerate}
        \item 证明:若函数 $g(x),f(x)$ 连续,则函数
              \[
                  \varphi(x)=\min\{f(x),g(x)\}, \qquad \psi(x)=\max\{f(x),g(x)\}
              \]
              连续;
        \item 设 $f_1(x),f_2(x),f_3(x)$ 在 $\lbrack a,b\rbrack$ 上连续,令函数 $f$ 的值等于三值 $f_1(x),f_2(x),f_3(x)$ 中介于其它两值之间的那个值,证明:$f$ 在 $\lbrack a,b\rbrack$ 上连续.
    \end{enumerate}
\end{example}
\exampleblank

\begin{solution}
    由恒等式
    \[
        \min\{f,g\}=\dfrac{f+g-|f-g|}{2},\qquad \max\{f,g\}=\dfrac{f+g+|f-g|}{2}
    \]
    可知两者连续.三个数的中间值可写为
    \[
        f_1+f_2+f_3-\max\{f_1,f_2,f_3\}-\min\{f_1,f_2,f_3\}.
    \]
    有限个连续函数的最大值和最小值仍连续,故所定义的 $f$ 连续.
\end{solution}

%例题3.1.8
\begin{example}
    是否存在 $\lbrack a,b\rbrack$ 上的连续函数:在有理点上取值为无理数,在无理点上取值为有理数?请作出判断并说明理由.
\end{example}
\exampleblank

\begin{solution}
    不存在.若 $f$ 连续且非常值,则其值域含有一个非退化区间.该区间中有不可数多个无理数 $y$,而每个这样的 $y$ 的原像都必须是有理数;这会从可数集 $\mathbb{Q}$ 向不可数集给出满射,矛盾.若 $f$ 为常值,又不可能同时在有理点取无理值,在无理点取有理值.
\end{solution}

%例题3.1.9
\begin{example}
    设 $f(x),g(x)$ 是定义在 $(a,b)$ 上的单调函数,且对 $\lbrack a,b\rbrack$ 中的任一有理数 $r$,均有 $f(r)=g(r)$,则 $f(x)$ 与 $g(x)$ 的连续点相同,在连续点上的函数值也相同.
\end{example}
\exampleblank

\begin{solution}
    对单调函数,任一点的左,右极限都可以分别用趋近该点的有理数列计算.由 $f(r)=g(r)$ 对所有有理数 $r$ 成立,得 $f(x-)=g(x-),f(x+)=g(x+)$.
    若 $f$ 在 $x$ 连续,两个单侧极限相等;单调性又迫使 $g(x)$ 等于该共同极限,故 $g$ 也连续且 $f(x)=g(x)$.交换 $f,g$ 即得逆命题.
\end{solution}

%例题3.1.10
\begin{example}
    设 $f(x)$ 定义在 $(a,b)$ 上,且对 $(a,b)$ 中任一点 $x_0$,均存在极限 $\displaystyle\lim\limits_{x\to x_0}f(x)$.若令
    \[
        g(x)=\lim\limits_{t\to x}f(t), \qquad x\in(a,b),
    \]
    则 $g(x)\in C((a,b))$.
\end{example}
\exampleblank

\begin{solution}
    记 $L=g(x_0)=\displaystyle\lim\limits_{t\to x_0}f(t)$.给定 $\varepsilon>0$,取 $\delta>0$,使 $0<|t-x_0|<2\delta$ 时 $|f(t)-L|<\varepsilon$.当 $|x-x_0|<\delta$ 时,令 $t\to x$ 且始终保持 $0<|t-x_0|<2\delta$,得 $|g(x)-L|\leq\varepsilon$.
    因而 $g(x)\to g(x_0)$,$g$ 在 $(a,b)$ 上连续.
\end{solution}



\subsection{确界函数}

%例题3.1.11
\begin{example}
    设 $f(x)$ 在 $\lbrack a,b\rbrack$ 上有界,对 $x\in(a,b\rbrack$ 定义
    \[
        M(x)=\sup\limits_{a\leq t<x}f(t),\qquad
        m(x)=\inf\limits_{a\leq t<x}f(t).
    \]
    则:
    \begin{enumerate}
        \item $M(x),m(x)$ 在 $(a,b\rbrack$ 上左连续;
        \item 若 $f(x)$ 在 $\lbrack a,b\rbrack$ 上连续,则 $M(x),m(x)$ 在 $(a,b)$ 上连续.
    \end{enumerate}
\end{example}
\exampleblank

\begin{solution}
    \leavevmode
    \begin{enumerate}
        \item 先证 $M$ 左连续.任取 $x_0\in(a,b\rbrack$ 及 $\varepsilon>0$.由上确界的定义,存在 $\xi\in\lbrack a,x_0)$,使
              \[
                  f(\xi)>M(x_0)-\varepsilon.
              \]
              取 $\delta=x_0-\xi>0$.当 $x\in(x_0-\delta,x_0)$ 时,有 $\xi<x<x_0$,于是
              \[
                  M(x_0)\geq M(x)\geq f(\xi)>M(x_0)-\varepsilon.
              \]
              因而 $|M(x)-M(x_0)|<\varepsilon$,故 $M$ 在 $x_0$ 左连续.由 $x_0$ 的任意性,$M$ 在 $(a,b\rbrack$ 上左连续.

              对 $-f$ 应用上述结论,注意
              \[
                  m(x)=-\sup\limits_{a\leq t<x}[-f(t)],
              \]
              即得 $m$ 在 $(a,b\rbrack$ 上左连续.

        \item 只需再证右连续.任取 $x_0\in(a,b)$.由于 $f$ 在 $\lbrack a,x_0\rbrack$ 上连续,
              \[
                  M(x_0)=\sup\limits_{a\leq t<x_0}f(t)
                  =\max\limits_{a\leq t\leq x_0}f(t)\geq f(x_0).
              \]
              给定 $\varepsilon>0$,由 $f$ 在 $x_0$ 连续,存在 $\delta>0$,使 $x_0+\delta<b$,且当 $t\in\lbrack x_0,x_0+\delta\rbrack$ 时,
              \[
                  f(t)<f(x_0)+\varepsilon\leq M(x_0)+\varepsilon.
              \]
              若 $x\in(x_0,x_0+\delta)$,则由 $M$ 单调不减,
              \[
                  M(x_0)\leq M(x)\leq M(x_0)+\varepsilon.
              \]
              故 $M$ 在 $x_0$ 右连续.结合第 (1) 问,$M$ 在 $(a,b)$ 上连续.同理对 $-f$ 应用该结论,得到 $m$ 在 $(a,b)$ 上连续.
    \end{enumerate}
\end{solution}

%例题3.1.12
\begin{example}
    设 $f(x)$ 在 $\lbrack 0,1\rbrack$ 上连续,且 $f(x)>0$,
    \[
        R(x)=\sup\limits_{0\leq y\leq x}f(y)\quad(0\leq x\leq1),
        \qquad
        G(x)=\lim\limits_{n\to\infty}\left[\dfrac{f(x)}{R(x)}\right]^n.
    \]
    试证:当且仅当 $f(x)$ 在 $\lbrack 0,1\rbrack$ 上单增时,$G$ 是连续的.
\end{example}
\exampleblank

\begin{solution}
    由 $f>0$ 且 $f(x)\leq R(x)$,
    \[
        G(x)= \begin{cases} 1, & f(x)=R(x), \\ 0, & f(x)<R(x). \end{cases}
    \]
    若 $f$ 单调不减,则 $R=f$,故 $G\equiv1$.反之,若 $G$ 连续,由区间的连通性及 $G$ 只取 $0,1$ 两值,$G$ 必为常值.又 $G(0)=1$,故 $G\equiv1$,从而 $f(x)=R(x)$ 对所有 $x$ 成立,这正表明 $f$ 单调不减.
\end{solution}

%例题3.1.13
\begin{example}
    已知
    \[
        f(x)=
        \begin{cases}
            x,   & 0\leq x<1,   \\
            k+1, & k\leq x<k+1,
        \end{cases}
        \qquad k=1,2,3,\ldots.
    \]
    求函数
    \[
        g(y)=\sup\limits_{f(x)\leq y}x
    \]
    在 $y\geq0$ 时的具体表达式,并指出 $g(y)$ 在各点处的左右连续性.
\end{example}
\exampleblank

\begin{solution}
    直接由定义分段讨论得
    \[
        g(y)= \begin{cases} y,                & 0\leq y<1, \\ \lfloor y\rfloor, & y\geq1. \end{cases}
    \]
    特别地,$g(1)=1$,故在 $y=1$ 连续.在每个整数 $m\geq2$ 处, $g(m-)=m-1,\,g(m)=g(m+)=m$,
    因而右连续但不左连续;其余点连续.
\end{solution}



\subsection{单射,介值性与单调性的关系}

\begin{proposition}[\markstar]
    设函数 $y=f(x)$ 在 $(a,b)$ 内有定义,且 $f$ 为单射并具有介值性,则:
    \begin{enumerate}
        \item $f(x)$ 严格单调,值域 $J=f((a,b))$ 为某个开区间;
        \item $f^{-1}(y)$ 在 $J$ 内严格单调,而且也具有介值性;
        \item $f(x)$ 与 $f^{-1}(y)$ 均连续.
    \end{enumerate}
\end{proposition}

\begin{corollary}[\markstar]
    若函数 $f(x)$ 在区间 $I$ 上处处连续且为单射,则 $f(x)$ 在 $I$ 上必严格单调.
\end{corollary}



\subsection{与点集结合 *}

%例题3.1.14
\begin{example}
    证明定理:$f(x)$ 在实轴 $x$ 上连续 $\Longleftrightarrow$ 任何开集的逆像仍为开集.
\end{example}
\exampleblank

\begin{solution}
    若 $f$ 连续,则对任意开集 $V$ 和 $x_0\in f^{-1}(V)$,可取 $f(x_0)$ 的小邻域包含在 $V$ 中,连续性给出 $x_0$ 的邻域包含在 $f^{-1}(V)$ 中,故逆像为开集.

    反之,若任一开集的逆像都开,则对任意 $\varepsilon>0$,区间 $(f(x_0)-\varepsilon,f(x_0)+\varepsilon)$ 的逆像是含 $x_0$ 的开集,故包含 $x_0$ 的某个邻域.这就是 $f$ 在 $x_0$ 连续.
\end{solution}

%例题3.1.15
\begin{example}
    设 $f(x)$ 在 $(-\infty,+\infty)$ 上有定义,证明:$f(x)$ 连续 $\Longleftrightarrow$ 对任意 $c\in(-\infty,+\infty)$,集合
    \[
        \{x\mid f(x)>c\} \quad\text{与}\quad \{x\mid f(x)<c\}
    \]
    为开集.
\end{example}
\exampleblank

\begin{solution}
    若 $f$ 连续,两个集合分别是开射线 $(c,+\infty)$ 与 $(-\infty,c)$ 的逆像,因而开.

    反之,由假设,任意开区间 $(\alpha,\beta)$ 的逆像
    \[
        f^{-1}((\alpha,\beta))=\{f>\alpha\}\cap\{f<\beta\}
    \]
    为开集.实轴上任意开集是开区间的并,故其逆像仍开,由上题得 $f$ 连续.
\end{solution}

%例题3.1.16
\begin{example}
    设 $f:\lbrack 0,1\rbrack\to\mathbb{R}$ 单调递增,且 $f(\lbrack 0,1\rbrack)$ 是区间,证明:$f$ 在 $\lbrack 0,1\rbrack$ 上连续.
\end{example}
\exampleblank

\begin{solution}
    反设 $f$ 在某点不连续.若 $x_0\in(0,1)$,由单调性知有限的单侧极限存在,且
    \[
        f(x_0-0)\leq f(x_0)\leq f(x_0+0).
    \]
    不连续意味着至少有一个不等号为严格不等号.若 $f(x_0-0)<f(x_0)$,则任取
    \[
        y\in(f(x_0-0),f(x_0)),
    \]
    由单调性可知 $x<x_0$ 时 $f(x)\leq f(x_0-0)<y$,而 $x\geq x_0$ 时 $f(x)\geq f(x_0)>y$,故 $y\notin f(\lbrack0,1\rbrack)$,这与值域为区间矛盾.情形 $f(x_0)<f(x_0+0)$ 同理.

    若间断发生在端点,例如 $x_0=0$,则必有 $f(0)<f(0+0)$,二者之间的值同样不可能属于值域;在 $x_0=1$ 时同理.故 $f$ 在 $\lbrack0,1\rbrack$ 上连续.
\end{solution}

%例题3.1.17
\begin{example}
    设 $y=f(x)$ 是从 $\lbrack 0,1\rbrack$ 到 $\lbrack 0,1\rbrack$ 的函数,其图像
    \[
        E=\{(x,y)\mid y=f(x),\ x\in\lbrack 0,1\rbrack\}
    \]
    是 $\lbrack 0,1\rbrack\times\lbrack 0,1\rbrack$ 上的闭集,求证:$f$ 是连续函数.
\end{example}
\exampleblank

\begin{solution}
    任取 $x_n\to x$.由 $f(x_n)\in\lbrack0,1\rbrack$ 的紧致性,任意子列都有收敛子子列 $f(x_{n_k})\to y$.此时 $(x_{n_k},f(x_{n_k}))\to(x,y)$.
    因图像 $E$ 闭,必有 $(x,y)\in E$,即 $y=f(x)$.因而 $f(x_n)$ 的每个收敛子列都只能收敛到 $f(x)$,故 $f(x_n)\to f(x)$.由序列判别法,$f$ 连续.
\end{solution}

%例题3.1.18
\begin{example}
    设 $y=f(x)$ 为 $X\to Y$ 的连续函数,$F$ 为 $Y$ 轴上的闭集,试证:$f^{-1}(F)$ 为 $X$ 轴上的闭集.
\end{example}
\exampleblank

\begin{solution}
    因 $F$ 闭,$Y\setminus F$ 开.连续性给出 $f^{-1}(Y\setminus F)$ 为开集,而 $X\setminus f^{-1}(F)=f^{-1}(Y\setminus F)$.
    故 $f^{-1}(F)$ 为闭集.
\end{solution}


\newpage

\section{连续性应用}


\subsection{\texorpdfstring{$f(f(x))$}{f(f(x))} 迭代问题}

\begin{proposition}[\markstar]
    设函数 $f,g$ 满足 $f(f(x))=g(x)$($f,g$ 定义合理),则:
    \begin{enumerate}
        \item 若 $g(x)$ 为单射,则 $f(x)$ 为单射;
        \item 若 $f(x)$ 为单调函数,则 $g(x)$ 一定为单调递增函数;
        \item 更一般地,若 $g(x)$ 为单射,且 $f(x)$ 具有介值性,则 $g(x)$ 一定为(严格)单调递增函数.
    \end{enumerate}
\end{proposition}

\begin{corollary}[\markstar]
    若 $g(x)$ 为单射,且 $f(x)$ 为连续函数,仍有 $g(x)$ 为(严格)单调递增函数.
\end{corollary}

%例题3.2.1
\begin{example}
    是否存在 $\mathbb{R}$ 上连续函数 $f(x)$,使得
    \[
        f(f(x))=e^{-x}.
    \]
\end{example}
\exampleblank

\begin{solution}
    不存在.因 $f\circ f=e^{-x}$ 是单射,$f$ 也必为单射.实轴上连续单射函数必严格单调;无论 $f$ 递增还是递减,$f\circ f$ 都递增,但 $e^{-x}$ 严格递减,矛盾.
\end{solution}

%例题3.2.2
\begin{example}
    设 $(-\infty,+\infty)$ 上的函数满足
    \[
        f(f(x))=-x,
    \]
    则 $f(x)$ 不是连续函数.
\end{example}
\exampleblank

\begin{solution}
    由 $f\circ f=-\operatorname{id}$ 为单射知 $f$ 单射.若 $f$ 连续,则 $f$ 严格单调,从而 $f\circ f$ 必严格递增;但 $-x$ 严格递减,矛盾.故 $f$ 不可能连续.
\end{solution}

%例题3.2.3
\begin{example}
    设 $f:\lbrack 0,1\rbrack\to\lbrack 0,1\rbrack$ 为连续函数,且 $f(0)=0,f(1)=1,f(f(x))\equiv x$,证明:$f(x)\equiv x$.
\end{example}
\exampleblank

\begin{solution}
    由 $f\circ f=\operatorname{id}$ 知 $f$ 为单射,故连续的 $f$ 严格单调.端点条件排除递减,因此 $f$ 严格递增.若某点 $f(x)>x$,对不等式作用递增函数 $f$ 得 $x=f(f(x))>f(x)$,矛盾;$f(x)<x$ 同理矛盾.故 $f(x)=x$.
\end{solution}

%例题3.2.4
\begin{example}
    设 $f$ 是从 $\mathbb{R}$ 到 $\mathbb{R}$ 的一对一连续映射,有不动点,且满足
    \[
        f(2x-f(x))\equiv x,\qquad \forall x\in\mathbb{R}.
    \]
    试证:$f(x)\equiv x\,(\forall x\in\mathbb{R})$.
\end{example}
\exampleblank

\begin{solution}
    由题设,对任意 $x\in\mathbb{R}$ 都有 $x=f(2x-f(x))$,故 $f$ 为满射.结合 $f$ 为一对一映射,$f$ 可逆,并且
    \[
        f^{-1}(x)=2x-f(x).
    \]
    连续单射函数 $f:\mathbb{R}\to\mathbb{R}$ 必严格单调.若 $f$ 严格递减,则 $f^{-1}$ 也严格递减,而 $2x-f(x)$ 严格递增,与上式矛盾.故 $f$ 严格递增.

    将题设中的 $x$ 换成 $f(x)$,由 $f^{-1}(f(x))=x$ 得
    \[
        2f(x)-f(f(x))=x,
    \]
    即
    \[
        f(f(x))=2f(x)-x.
    \]
    令 $d(x)=f(x)-x$,则
    \[
        d(f(x))=f(f(x))-f(x)=f(x)-x=d(x).
    \]
    又由 $f^{-1}(x)=2x-f(x)=x-d(x)$ 可知 $d(f^{-1}(x))=d(x)$.因此对任意整数 $n\in\mathbb{Z}$,归纳可得
    \[
        f^n(x)=x+nd(x),
    \]
    其中 $n<0$ 时 $f^n$ 表示 $f^{-1}$ 的迭代.

    设 $c$ 为 $f$ 的一个不动点.由于 $f$ 及 $f^{-1}$ 均严格递增,对任意 $n\in\mathbb{Z}$,$f^n$ 也严格递增且 $f^n(c)=c$.因此若 $x>c$,则 $f^n(x)>c$ 对所有 $n\in\mathbb{Z}$ 成立;若 $x<c$,则 $f^n(x)<c$ 对所有 $n\in\mathbb{Z}$ 成立.

    若存在 $x$ 使 $d(x)\neq0$,则数列 $x+nd(x)$ 随整数 $n$ 可取到 $c$ 的两侧,这与上面的保序性矛盾.故对任意 $x\in\mathbb{R}$ 都有 $d(x)=0$,即
    \[
        f(x)\equiv x.
    \]
\end{solution}

%例题3.2.5
\begin{example}
    设 $f\in C(( -\infty,+\infty))$ 且
    \[
        f(f(f(x)))=x, \qquad x\in(-\infty,+\infty),
    \]
    则 $f(x)\equiv x$.
\end{example}
\exampleblank

\begin{solution}
    由 $f^3=\operatorname{id}$ 知 $f$ 单射,因而严格单调.若 $f$ 递减,则 $f^3$ 递减,与恒等函数递增矛盾;故 $f$ 递增.若 $f(x)>x$,则 $x<f(x)<f^2(x)<f^3(x)=x$,
    矛盾;$f(x)<x$ 同理.故 $f(x)=x$.
\end{solution}

%例题3.2.6
\begin{example}
    设 $f:\lbrack 0,1\rbrack\to\lbrack 0,1\rbrack$ 连续,且 $f(0)=0,f(1)=1$,定义
    \[
        f^n(x)\overset{\triangle}{=}(f\circ f\circ\cdots\circ f)(x) \quad(n\text{ 次复合}).
    \]
    若存在 $n_0$,使得 $f^{n_0}(x)=x\,(0\leq x\leq1)$,则 $f(x)=x$.
\end{example}
\exampleblank

\begin{solution}
    由 $f^{n_0}=\operatorname{id}$ 知 $f$ 单射,故 $f$ 严格单调.由 $f(0)=0,f(1)=1$ 知 $f$ 递增.若 $f(x)>x$,则递增性给出 $x<f(x)<f^2(x)<\cdots<f^{n_0}(x)=x$,
    矛盾;$f(x)<x$ 同理.因而 $f(x)=x$.
\end{solution}

%例题3.2.7
\begin{example}
    设 $f:\mathbb{R}\to\mathbb{R}$ 连续,且 $f(x)$ 无不动点,
    \[
        f^n(x)\overset{\triangle}{=}(f\circ f\circ\cdots\circ f)(x),
    \]
    则对任意 $x_0\in\mathbb{R}$,数列 $x_n=f^n(x_0)$ 无界.
\end{example}
\exampleblank

\begin{solution}
    连续函数 $f(x)-x$ 无零点,故在连通集 $\mathbb{R}$ 上恒正或恒负.因而迭代列 $x_{n+1}=f(x_n)$ 严格单调.若它有界,则存在有限极限 $L$,由连续性
    \[
        L=\lim\limits_{n\to\infty}x_{n+1} =\lim\limits_{n\to\infty}f(x_n)=f(L),
    \]
    这与无不动点矛盾.故迭代列无界.
\end{solution}

%例题3.2.8
\begin{example}
    设 $f(x)$ 是 $\mathbb{R}$ 上连续函数,且
    \[
        \lim\limits_{x\to+\infty}f(f(x))=+\infty.
    \]
    证明:
    \[
        \lim\limits_{x\to+\infty}|f(x)|=+\infty.
    \]
\end{example}
\exampleblank

\begin{solution}
    反设 $\displaystyle\lim\limits_{x\to+\infty}|f(x)|\neq+\infty$.则存在常数 $M>0$ 及数列 $x_n\to+\infty$,使得
    \[
        |f(x_n)|\leq M.
    \]
    因此数列 $\{f(x_n)\}$ 有界.由 Bolzano--Weierstrass 定理,可取子列 $\{x_{n_k}\}$,使
    \[
        f(x_{n_k})\to A
    \]
    对某个 $A\in\mathbb{R}$ 成立.由 $f$ 的连续性,
    \[
        f(f(x_{n_k}))\to f(A),
    \]
    右端为有限数.另一方面,由于 $x_{n_k}\to+\infty$,题设给出
    \[
        f(f(x_{n_k}))\to+\infty,
    \]
    矛盾.故
    \[
        \lim\limits_{x\to+\infty}|f(x)|=+\infty.
    \]
\end{solution}


\subsection{周期函数}

\begin{proposition}[\markstar]
    非常数连续周期函数必有最小正周期.
\end{proposition}

\begin{corollary}[\markstar]
    \leavevmode
    \begin{enumerate}
        \item 若 $f$(非常数)是连续函数,有最小正周期 $T_0$,则对任意 $T\in\{f\text{ 的正周期}\}$,必存在 $n\in\mathbb{N}$,使得 $T=nT_0$;
        \item 设 $f$(非常数)是连续函数,有周期 $T>0$,若 $\alpha T$ 也是 $f$ 的周期,则 $\alpha\in\mathbb{Q}$.
    \end{enumerate}
\end{corollary}

%例题3.2.9
\begin{example}
    证明命题 3.2.2 及推论 3.2.2.
\end{example}
\exampleblank

\begin{solution}
    设 $\mathcal P$ 为 $f$ 的正周期全体,$T_0=\inf\mathcal P$.

    先证 $T_0>0$.若 $T_0=0$,则可取 $T_n\in\mathcal P$,使
    \[
        0<T_n<\dfrac{1}{n},
    \]
    从而 $T_n\to0$.任取 $x,y\in\mathbb{R}$,令
    \[
        m_n=\left\lfloor\dfrac{y-x}{T_n}\right\rfloor\in\mathbb{Z}.
    \]
    由取整函数的定义,
    \[
        0\leq\dfrac{y-x}{T_n}-m_n<1,
    \]
    即
    \[
        0\leq y-(x+m_nT_n)<T_n.
    \]
    因此 $x+m_nT_n\to y$.由于整数倍周期仍为周期,
    \[
        f(x+m_nT_n)=f(x).
    \]
    令 $n\to\infty$,由连续性得 $f(y)=f(x)$.因 $x,y$ 任意,$f$ 为常数,与命题中 $f$ 非常数矛盾.故 $T_0>0$.

    由下确界的定义,可取 $T_n\in\mathcal P$,使
    \[
        T_0\leq T_n<T_0+\dfrac{1}{n}.
    \]
    于是 $T_n\to T_0$.对任意 $x\in\mathbb{R}$,
    \[
        f(x+T_0)=\lim\limits_{n\to\infty}f(x+T_n)=f(x),
    \]
    故 $T_0$ 本身也是周期,从而是最小正周期.

    再证推论.任取正周期 $T$,作带余除法
    \[
        T=qT_0+r,\qquad q=\left\lfloor\dfrac{T}{T_0}\right\rfloor,\qquad 0\leq r<T_0.
    \]
    若 $r>0$,由于 $T$ 与 $qT_0$ 都是周期,其差 $r$ 也是正周期,这与 $T_0$ 的最小性矛盾.故 $r=0$,即
    \[
        T=qT_0,\qquad q\in\mathbb{N}.
    \]
    因此若 $T$ 与 $\alpha T$ 都是正周期,则存在 $m,n\in\mathbb{N}$,使
    \[
        T=mT_0,\qquad \alpha T=nT_0,
    \]
    从而
    \[
        \alpha=\dfrac{n}{m}\in\mathbb{Q}.
    \]
\end{solution}

%例题3.2.10
\begin{example}
    试证明下列命题:
    \begin{enumerate}
        \item 设 $f\in C((-\infty,+\infty))$,若 $T_1,T_2$ 皆为 $f$ 的周期,且 $T_1$ 与 $T_2$ 不可公约(即 $\dfrac{T_1}{T_2}\notin\mathbb{Q}$),则 $f(x)$ 为常数;
        \item 设 $f\in C((-\infty,+\infty))$ 是周期函数,其正周期全体形成的数集的下确界为 $T_0$.若 $T_0=0$,则 $f(x)\equiv C$(常数).
    \end{enumerate}
\end{example}
\exampleblank

\begin{solution}
    \leavevmode
    \begin{enumerate}
        \item 反设 $f$ 非常数.由命题 3.2.2,连续非常数周期函数存在最小正周期 $T_0>0$.再由推论 3.2.2,任意正周期都是 $T_0$ 的正整数倍,故存在 $m,n\in\mathbb{N}$,使
              \[
                  T_1=mT_0,\qquad T_2=nT_0.
              \]
              因此
              \[
                  \dfrac{T_1}{T_2}=\dfrac{m}{n}\in\mathbb{Q},
              \]
              与题设矛盾.故 $f$ 只能为常数.

        \item 反设 $f$ 非常数.仍由命题 3.2.2,$f$ 存在最小正周期 $T_*>0$.于是任意正周期 $T$ 都满足 $T\geq T_*$,从而正周期集的下确界满足
              \[
                  T_0\geq T_*>0,
              \]
              与题设 $T_0=0$ 矛盾.故 $f$ 为常数.
    \end{enumerate}
\end{solution}

%例题3.2.11
\begin{example}
    设 $f(x),g(x)$ 均为非常数连续周期函数,其最小正周期分别为 $T_1,T_2$.若
    \[
        \dfrac{T_1}{T_2}\notin\mathbb{Q},
    \]
    证明:$h(x)=f(x)+g(x)$ 不是周期函数.
\end{example}
\exampleblank

\begin{solution}
    反设 $h$ 有正周期 $T$.由 $h(x+T)=h(x)$,
    \[
        f(x+T)-f(x)=-[g(x+T)-g(x)].
    \]
    记
    \[
        F(x)=f(x+T)-f(x).
    \]
    由于 $T_1$ 是 $f$ 的周期,$F$ 以 $T_1$ 为周期;另一方面由上式及 $T_2$ 是 $g$ 的周期,$F$ 也以 $T_2$ 为周期.由例题 3.2.10(1) 及 $\dfrac{T_1}{T_2}\notin\mathbb{Q}$,连续函数 $F$ 必为常数,设 $F(x)\equiv c$.

    于是对任意正整数 $n$,
    \[
        f(x+nT)=f(x)+nc.
    \]
    连续周期函数 $f$ 有界,故只能有 $c=0$.因此 $T$ 是 $f$ 的周期.由
    \[
        g(x+T)-g(x)=-F(x)=0
    \]
    可知 $T$ 也是 $g$ 的周期.

    再由推论 3.2.2,存在 $m,n\in\mathbb{N}$,使
    \[
        T=mT_1=nT_2,
    \]
    从而 $\dfrac{T_1}{T_2}=\dfrac{n}{m}\in\mathbb{Q}$,矛盾.故 $h$ 不是周期函数.
\end{solution}

\subsection{Cauchy 方程}

\begin{proposition}[\markstar]
    设 $f(x)$ 定义在 $(-\infty,+\infty)$ 上满足 Cauchy 方程
    \[
        f(x+y)=f(x)+f(y), \qquad x,y\in(-\infty,+\infty).
    \]
    则 $f(x)=f(1)x$ 对任意 $x\in\mathbb{Q}$ 成立.进一步,若 $f$ 还满足如下条件之一:
    \begin{enumerate}
        \item $f$ 在某点 $x=x_0$ 处连续;
        \item $f(x)$ 在某个区间上有界;
        \item $f(x)$ 单调,
    \end{enumerate}
    则 $f(x)=f(1)x$ 对任意 $x\in\mathbb{R}$ 成立.
\end{proposition}

\begin{remark}
    事实上由于 Cauchy 方程,$f$ 满足平移性,故 $f$ 在某点的性质可以推广到整个 $\mathbb{R}$ 上,如 $f$ 在 $x=0$ 处连续 $\Longleftrightarrow$ $f$ 在 $\mathbb{R}$ 上处处连续,本质上十分类似后续课程中的线性泛函.
\end{remark}

%例题3.2.12
\begin{example}
    设 $f:\mathbb{R}\to\mathbb{R}$ 满足
    \[
        f(x+y)=f(x)+f(y),\qquad \forall x,y\in\mathbb{R},
    \]
    且 $f$ 在 $x=0$ 处连续.证明:$f$ 在 $\mathbb{R}$ 上连续,且
    \[
        f(x)=f(1)x,\qquad \forall x\in\mathbb{R}.
    \]
\end{example}
\exampleblank

\begin{solution}
    令 $x=y=0$,得 $f(0)=0$.对任意固定的 $x\in\mathbb{R}$,
    \[
        f(x+h)-f(x)=f(h).
    \]
    当 $h\to0$ 时,由 $f$ 在 $0$ 处连续及 $f(0)=0$,
    \[
        f(h)\to0.
    \]
    因而 $f(x+h)\to f(x)$,故 $f$ 在 $\mathbb{R}$ 上处处连续.

    由 Cauchy 方程归纳可得
    \[
        f(nx)=nf(x),\qquad n\in\mathbb{Z}.
    \]
    取 $x=1$,对任意有理数 $r=\dfrac{p}{q}$($p\in\mathbb{Z},q\in\mathbb{N}$),有
    \[
        qf(r)=f(qr)=f(p)=pf(1),
    \]
    从而
    \[
        f(r)=rf(1).
    \]
    任取 $x\in\mathbb{R}$,取有理数列 $r_n\to x$.由连续性,
    \[
        f(x)=\lim\limits_{n\to\infty}f(r_n)
        =\lim\limits_{n\to\infty}r_nf(1)
        =xf(1).
    \]
\end{solution}

%例题3.2.13
\begin{example}
    设 $f\in C((-\infty,+\infty))$,且
    \[
        f\left(\dfrac{x+y}{2}\right)=\dfrac{1}{2}[f(x)+f(y)],\qquad \forall x,y\in\mathbb{R},
    \]
    证明:
    \[
        f(x)=[f(1)-f(0)]x+f(0).
    \]
\end{example}
\exampleblank

\begin{solution}
    对任意 $x,y\in\mathbb{R}$,由题设
    \[
        f\left(\dfrac{x+y}{2}\right)=\dfrac{f(x)+f(y)}{2}.
    \]
    另一方面,把题设中的两个自变量取为 $x+y$ 与 $0$,又有
    \[
        f\left(\dfrac{x+y}{2}\right)=\dfrac{f(x+y)+f(0)}{2}.
    \]
    比较两式得
    \[
        f(x+y)+f(0)=f(x)+f(y).
    \]
    令
    \[
        g(x)=f(x)-f(0),
    \]
    则 $g$ 连续且满足
    \[
        g(x+y)=g(x)+g(y).
    \]
    由例题 3.2.12,
    \[
        g(x)=g(1)x=[f(1)-f(0)]x.
    \]
    因此
    \[
        f(x)=[f(1)-f(0)]x+f(0).
    \]
\end{solution}

%例题3.2.14
\begin{example}
    设 $f\in C((0,+\infty))$,且
    \[
        f(xy)=f(x)+f(y),\qquad \forall x,y>0.
    \]
    证明:存在常数 $c\in\mathbb{R}$,使得
    \[
        f(x)=c\ln x,\qquad x>0.
    \]
\end{example}
\exampleblank

\begin{solution}
    令
    \[
        g(t)=f(e^t),\qquad t\in\mathbb{R}.
    \]
    由于 $f$ 连续,$g$ 也连续.并且对任意 $s,t\in\mathbb{R}$,
    \[
        g(s+t)=f(e^{s+t})=f(e^se^t)=f(e^s)+f(e^t)=g(s)+g(t).
    \]
    由例题 3.2.12,存在常数 $c=g(1)$,使
    \[
        g(t)=ct,\qquad t\in\mathbb{R}.
    \]
    对任意 $x>0$,取 $t=\ln x$,则
    \[
        f(x)=f(e^t)=g(t)=ct=c\ln x.
    \]
    反之,$f(x)=c\ln x$ 显然满足题设方程,故这就是全部连续解.
\end{solution}

%例题3.2.15
\begin{example}
    求函数方程
    \[
        f(x+y)+f(x-y)=2f(x)f(y),\qquad \forall x,y\in\mathbb{R},
    \]
    在实轴上的非零连续解.
\end{example}
\exampleblank

\begin{solution}
    设 $f$ 为非零连续解.在方程中令 $y=0$,得
    \[
        2f(x)=2f(x)f(0).
    \]
    由于 $f$ 非零,故 $f(0)=1$.再令 $x=0$,得到
    \[
        f(y)+f(-y)=2f(y),
    \]
    因而 $f(-y)=f(y)$,$f$ 为偶函数.

    由于 $f(0)=1$ 且 $f$ 连续,可取 $a>0$,使
    \[
        A=\int_0^a f(t)\,\mathrm{d}t\neq0.
    \]
    对原方程关于 $y$ 在 $\lbrack0,a\rbrack$ 上积分,得
    \[
        \int_0^a[f(x+y)+f(x-y)]\,\mathrm{d}y
        =2f(x)\int_0^a f(y)\,\mathrm{d}y.
    \]
    左端合并为 $\displaystyle\int_{x-a}^{x+a}f(t)\,\mathrm{d}t$,故
    \[
        f(x)=\dfrac{1}{2A}\int_{x-a}^{x+a}f(t)\,\mathrm{d}t.
    \]
    由此
    \[
        f'(x)=\dfrac{f(x+a)-f(x-a)}{2A},
    \]
    所以 $f\in C^1(\mathbb{R})$.再对上式求导可知 $f\in C^2(\mathbb{R})$.

    于是可对原方程关于 $y$ 求两次导数:
    \[
        f''(x+y)+f''(x-y)=2f(x)f''(y).
    \]
    令 $y=0$,得到
    \[
        f''(x)=cf(x),\qquad c=f''(0).
    \]
    又因为 $f$ 为偶函数,
    \[
        f(0)=1,\qquad f'(0)=0.
    \]

    若 $c=-\lambda^2<0$,上述初值问题的唯一解为
    \[
        f(x)=\cos(\lambda x).
    \]
    若 $c=0$,则 $f(x)\equiv1$,这就是 $\lambda=0$ 时的 $\cos(\lambda x)$.若 $c=\lambda^2>0$,则
    \[
        f(x)=\cosh(\lambda x).
    \]
    这些函数由相应的加法公式均满足原方程.因此全部非零连续解为
    \[
        f(x)=\cos(ax)\quad\text{或}\quad f(x)=\cosh(ax),\qquad a\in\mathbb{R}.
    \]
\end{solution}

\newpage

\section{闭区间上连续函数性质}


\subsection{有界性,最值性}

%例题3.3.1
\begin{example}
    设 $f\in C(\lbrack a,b\rbrack)$,若对任意 $x\in\lbrack a,b\rbrack$,存在 $x'\in\lbrack a,b\rbrack$,使得
    \[
        |f(x')|\leq\dfrac{|f(x)|}{2},
    \]
    则存在 $x_0\in\lbrack a,b\rbrack$,使得 $f(x_0)=0$.
\end{example}
\exampleblank

\begin{solution}
    由连续性,$|f|$ 在紧区间上取到最小值 $m\geq0$.在最小点 $x$ 应用题设,存在 $x'$ 使 $m\leq |f(x')|\leq\dfrac{m}{2}$.
    因而 $m=0$,最小点即为所求零点.
\end{solution}

%例题3.3.2
\begin{example}
    设 $f\in C(\lbrack a,b\rbrack)$,若 $(a,b)$ 中无 $f(x)$ 的极值点,则 $f(x)$ 在 $\lbrack a,b\rbrack$ 上单调.
\end{example}
\exampleblank

\begin{solution}
    反设 $f$ 不单调,则存在 $x_1<x_2<x_3$ 使 $f(x_2)$ 严格大于两端值或严格小于两端值.连续函数在 $\lbrack x_1,x_3\rbrack$ 上取到最值或最小值,且该极值不可能只在端点取到,故 $(x_1,x_3)$ 内存在极值点,矛盾.因而 $f$ 单调.
\end{solution}

%例题3.3.3
\begin{example}
    设 $f\in C(( -\infty,+\infty))$,若对任意的开区间 $(a,b)$,其值域 $f((a,b))$ 必是开区间,则 $f(x)$ 必单调.
\end{example}
\exampleblank

\begin{solution}
    若 $f$ 不单调,则存在三点 $x_1<x_2<x_3$,使中间值严格高于或低于两端值.取包含 $x_2$ 的足够小开区间 $I$,$f$ 在 $I$ 内取到一个局部最大或最小值,于是 $f(I)$ 含有自己的一个端点,不可能是开区间.矛盾,故 $f$ 必单调.
\end{solution}



\subsection{介值性,零点存在性}

%例题3.3.4
\begin{example}
    \begin{enumerate}
        \item 设 $f(x)$ 在 $\lbrack a,b\rbrack$ 上单调,若 $f(x)$ 在 $\lbrack a,b\rbrack$ 上具有介值性,则 $f\in C(\lbrack a,b\rbrack)$;
        \item 设 $f(x)$ 定义在 $\lbrack a,b\rbrack$ 上,若对任意 $x\in(a,b)$,均存在极限 $\displaystyle\lim\limits_{t\to x}f(t)$,且又具有介值性,则 $f\in C((a,b))$.
    \end{enumerate}
\end{example}
\exampleblank

\begin{solution}
    \leavevmode
    \begin{enumerate}
        \item 单调函数若在某点不连续,则该点必为跳跃间断点,从而会漏掉左,右极限之间的一整段函数值,这与介值性矛盾.故 $f$ 在 $\lbrack a,b\rbrack$ 上连续.

        \item 任取 $x_0\in(a,b)$,记
              \[
                  L=\lim\limits_{t\to x_0}f(t).
              \]
              只需证明 $L=f(x_0)$.反设 $L\neq f(x_0)$.不妨设 $L<f(x_0)$,取
              \[
                  c\in(L,f(x_0)).
              \]
              令 $\varepsilon=c-L>0$.由极限的定义,存在 $\delta>0$,使 $0<|t-x_0|<\delta$ 时
              \[
                  |f(t)-L|<\varepsilon,
              \]
              从而 $f(t)<c$.

              取任意 $t\in(x_0,x_0+\delta)$.由于 $f(t)<c<f(x_0)$,而 $f$ 具有介值性,在区间 $\lbrack x_0,t\rbrack$ 内应存在 $\xi\in(x_0,t)$,使 $f(\xi)=c$.但 $0<|\xi-x_0|<\delta$,按上面的极限估计又有 $f(\xi)<c$,矛盾.情形 $L>f(x_0)$ 同理.

              因此 $L=f(x_0)$,$f$ 在 $x_0$ 连续.由 $x_0$ 的任意性,$f\in C((a,b))$.
    \end{enumerate}
\end{solution}

%例题3.3.5
\begin{example}
    设 $f\in C(\lbrack a,b\rbrack)$,$x_i\in\lbrack a,b\rbrack\,(i=1,2,\ldots,n)$,则存在 $\xi\in\lbrack a,b\rbrack$,使得
    \[
        f(\xi)=\dfrac{1}{n}\sum\limits_{i=1}^n f(x_i).
    \]
\end{example}
\exampleblank

\begin{solution}
    记 $m=\min f$,$M=\max f$.每个 $f(x_i)$ 都在 $\lbrack m,M\rbrack$ 中,故其算术平均值 $A$ 也在 $\lbrack m,M\rbrack$ 中.由连续函数的介值定理,存在 $\xi\in\lbrack a,b\rbrack$ 使 $f(\xi)=A$.
\end{solution}

%例题3.3.6
\begin{example}
    \begin{enumerate}
        \item 设 $f\in C(\lbrack 0,1\rbrack)$,若 $f(0)=f(1)$,则存在 $\alpha,\beta$ 满足
              \[
                  0\leq\alpha<\beta\leq1, \qquad \beta-\alpha=\dfrac{1}{5},
              \]
              使得 $f(\alpha)=f(\beta)$;
        \item 设 $f(x)$ 在 $(-\infty,+\infty)$ 上以 $T$ 为周期的连续函数,则存在 $\xi\in(-\infty,+\infty)$ 使得
              \[
                  f\left(\xi+\dfrac{T}{2}\right)=f(\xi).
              \]
    \end{enumerate}
\end{example}
\exampleblank

\begin{solution}
    第 (1) 问令 $h(x)=f\left(x+\dfrac{1}{5}\right)-f(x)$,$0\leq x\leq\dfrac{4}{5}$.有
    \[
        \displaystyle\sum\limits_{k=0}^{4}h\left(\dfrac{k}{5}\right)=f(1)-f(0)=0.
    \]
    因而这些值中要么有零,要么既有正值又有负值;由 $h$ 连续必有零点.

    第 (2) 问令 $h(x)=f\left(x+\dfrac{T}{2}\right)-f(x)$.周期性给出 $h\left(x+\dfrac{T}{2}\right)=-h(x)$,故由介值定理,$h$ 在 $x$ 与 $x+\dfrac{T}{2}$ 之间必有零点.
\end{solution}

%例题3.3.7
\begin{example}
    设 $f\in C(\lbrack 0,1\rbrack)$,$n\in\mathbb{N}$.
    \begin{enumerate}
        \item 若 $f(0)=f(1)$,则存在 $\xi_n\in\left[0,1-\dfrac{1}{n}\right]$,使得
              \[
                  f\left(\xi_n+\dfrac{1}{n}\right)=f(\xi_n);
              \]
        \item 若 $f(0)=0,f(1)=1$,则存在 $\xi_n\in\left[0,1-\dfrac{1}{n}\right]$,使得
              \[
                  f\left(\xi_n+\dfrac{1}{n}\right)=f(\xi_n)+\dfrac{1}{n}.
              \]
    \end{enumerate}
\end{example}
\exampleblank

\begin{solution}
    \leavevmode
    \begin{enumerate}
        \item 令
              \[
                  h(x)=f\left(x+\dfrac{1}{n}\right)-f(x),
                  \qquad
                  0\leq x\leq1-\dfrac{1}{n}.
              \]
              在分点 $0,\dfrac{1}{n},\ldots,\dfrac{n-1}{n}$ 处有望远求和
              \[
                  \sum\limits_{k=0}^{n-1}h\left(\dfrac{k}{n}\right)
                  =\sum\limits_{k=0}^{n-1}
                  \left[
                      f\left(\dfrac{k+1}{n}\right)
                      -f\left(\dfrac{k}{n}\right)
                  \right]
                  =f(1)-f(0)=0.
              \]
              若其中某一项为 $0$,结论立即成立.否则这些有限个数中必有正值和负值.由于 $h$ 连续,由介值定理存在
              \[
                  \xi_n\in\left[0,1-\dfrac{1}{n}\right]
              \]
              使 $h(\xi_n)=0$,即
              \[
                  f\left(\xi_n+\dfrac{1}{n}\right)=f(\xi_n).
              \]

        \item 令
              \[
                  g(x)=f(x)-x.
              \]
              则 $g\in C(\lbrack0,1\rbrack)$ 且 $g(0)=g(1)=0$.对 $g$ 应用第 (1) 问,存在
              \[
                  \xi_n\in\left[0,1-\dfrac{1}{n}\right]
              \]
              使
              \[
                  g\left(\xi_n+\dfrac{1}{n}\right)=g(\xi_n).
              \]
              展开即得
              \[
                  f\left(\xi_n+\dfrac{1}{n}\right)
                  =f(\xi_n)+\dfrac{1}{n}.
              \]
    \end{enumerate}
\end{solution}

%例题3.3.8
\begin{example}
    设 $f\in C(( -\infty,+\infty))$,且以 $1$ 为周期,则存在 $\xi$,使得
    \[
        f(\xi+\pi)=f(\xi).
    \]
\end{example}
\exampleblank

\begin{solution}
    令 $h(x)=f(x+\pi)-f(x)$.由 $f$ 为 $1$ 周期函数,$h$ 也为 $1$ 周期函数,且
    \[
        \displaystyle\int_0^1h(x)\,\mathrm{d}x
        =\int_0^1f(x+\pi)\,\mathrm{d}x-\int_0^1f(x)\,\mathrm{d}x=0.
    \]
    连续函数 $h$ 若无零点,则在 $\lbrack0,1\rbrack$ 上恒正或恒负,与积分为零矛盾.故存在 $\xi$ 使 $h(\xi)=0$.
\end{solution}

%例题3.3.9
\begin{example}
    设 $f,g\in C(\lbrack 0,1\rbrack)$,且
    \[
        \max\limits_{x\in\lbrack 0,1\rbrack}f(x)=\max\limits_{x\in\lbrack 0,1\rbrack}g(x),
    \]
    试证:存在 $x_0\in\lbrack 0,1\rbrack$ 使得
    \[
        e^{f(x_0)}+3f(x_0)=e^{g(x_0)}+3g(x_0).
    \]
\end{example}
\exampleblank

\begin{solution}
    记公共最大值为 $M$,并令 $\Phi(t)=e^t+3t$.因 $\Phi'(t)=e^t+3>0$,$\Phi$ 严格递增.若 $a,b$ 分别为 $f,g$ 的最大值点,则
    \[
        \Phi(f(a))-\Phi(g(a))\geq0, \qquad \Phi(f(b))-\Phi(g(b))\leq0.
    \]
    对连续函数 $\Phi(f)-\Phi(g)$ 使用介值定理,即得所求 $x_0$.
\end{solution}

%例题3.3.10
\begin{example}
    设函数 $f(x)$ 在 $(-\infty,+\infty)$ 上连续,$n$ 为奇数,若
    \[
        \displaystyle\lim\limits_{x\to+\infty}\dfrac{f(x)}{x^n}
        =\lim\limits_{x\to-\infty}\dfrac{f(x)}{x^n}=1,
    \]
    试证:方程 $f(x)+x^n=0$ 有实根.
\end{example}
\exampleblank

\begin{solution}
    令 $F(x)=f(x)+x^n$.由题设及 $n$ 为奇数,
    \[
        \displaystyle\lim\limits_{x\to+\infty}\dfrac{F(x)}{x^n}=2,
        \qquad
        \lim\limits_{x\to-\infty}\dfrac{F(x)}{x^n}=2.
    \]
    因而 $F(x)\to+\infty$ $(x\to+\infty)$,而 $F(x)\to-\infty$ $(x\to-\infty)$.由介值定理,$F$ 必有零点.
\end{solution}

%例题3.3.11
\begin{example}
    设
    \[
        f_n(x)=x+x^2+\cdots+x^n, \qquad n=2,3,\ldots.
    \]
    \begin{enumerate}
        \item 证明:方程 $f_n(x)=1$ 在 $\lbrack 0,+\infty)$ 上有唯一的实根 $x_n$;
        \item 证明数列 $\{x_n\}$ 有极限,并求出 $\displaystyle\lim\limits_{n\to+\infty}x_n$.
    \end{enumerate}
\end{example}
\exampleblank

\begin{solution}
    $f_n$ 在 $\lbrack0,+\infty)$ 上严格递增,且 $f_n(0)=0<1<f_n(1)=n$,故有唯一根 $x_n\in(0,1)$.又 $f_{n+1}(x_n)=1+x_n^{n+1}>1$,
    故 $x_{n+1}<x_n$,从而 $x_n\to L\in\lbrack0,1)$.由 $1=\dfrac{x_n(1-x_n^n)}{1-x_n}$
    令 $n\to\infty$,得 $1=\dfrac{L}{1-L}$,故 $L=\dfrac{1}{2}$.
\end{solution}

%例题3.3.12
\begin{example}
    设
    \[
        f_n(x)=\cos x+\cos^2x+\cdots+\cos^n x,
    \]
    求证:
    \begin{enumerate}
        \item 对任意 $n$,$f_n(x)=1$ 在 $\left[0,\dfrac{\pi}{3}\right)$ 内有且仅有一根;
        \item 若 $x_n\in\left[0,\dfrac{\pi}{3}\right)$ 是 $f_n(x)=1$ 的根,则
              \[
                  \displaystyle\lim\limits_{n\to\infty}x_n=\dfrac{\pi}{3}.
              \]
    \end{enumerate}
\end{example}
\exampleblank

\begin{solution}
    令 $y=\cos x$.在 $\lbrack0,\dfrac{\pi}{3})$ 上,$y$ 由 $1$ 严格减到 $\dfrac{1}{2}$,方程化为 $y+y^2+\cdots+y^n=1$.
    左边对 $y>0$ 严格递增,在 $y=\dfrac{1}{2}$ 时小于 $1$,在 $y=1$ 时大于 $1$,故唯一根 $y_n\in\left(\dfrac{1}{2},1\right)$.由上题同样的论证,$y_n\to\dfrac{1}{2}$,因此 $x_n=\arccos y_n\to\dfrac{\pi}{3}$.
\end{solution}

%例题3.3.13
\begin{example}
    设 $f(x)$ 在 $(0,+\infty)$ 连续有界,试证:对任意 $T\in\mathbb{R}$,存在 $x_n\to+\infty$,使得
    \[
        \lim\limits_{n\to\infty}[f(x_n+T)-f(x_n)]=0.
    \]
\end{example}
\exampleblank

\begin{solution}
    当 $T=0$ 时结论显然.先设 $T>0$,令
    \[
        g(x)=f(x+T)-f(x),\qquad x>0.
    \]
    反设不存在满足题意的数列.则不存在趋于 $+\infty$ 的数列 $\{x_n\}$ 使 $g(x_n)\to0$,因此存在 $\varepsilon_0>0$ 和 $A>0$,使得
    \[
        |g(x)|\geq\varepsilon_0,\qquad x\geq A.
    \]
    事实上,若对任意 $m\in\mathbb{N}$ 都存在 $x_m\geq m$ 使 $|g(x_m)|<\dfrac{1}{m}$,则 $x_m\to+\infty$ 且 $g(x_m)\to0$,与反设矛盾.

    函数 $g$ 在连通区间 $\lbrack A,+\infty)$ 上连续且不为 $0$,故其符号恒定.若 $g(x)\geq\varepsilon_0$,则对任意正整数 $m$,
    \[
        \begin{aligned}
            f(A+mT)-f(A)
            &=\sum\limits_{k=0}^{m-1}
            [f(A+(k+1)T)-f(A+kT)]\\
            &=\sum\limits_{k=0}^{m-1}g(A+kT)
            \geq m\varepsilon_0.
        \end{aligned}
    \]
    这与 $f$ 有界矛盾.若 $g(x)\leq-\varepsilon_0$,同理得到
    \[
        f(A+mT)-f(A)\leq-m\varepsilon_0,
    \]
    仍与有界性矛盾.故 $T>0$ 时结论成立.

    若 $T<0$,令 $S=-T>0$.由已证结论,存在 $y_n\to+\infty$,使
    \[
        f(y_n+S)-f(y_n)\to0.
    \]
    令 $x_n=y_n+S$,则 $x_n\to+\infty$,且
    \[
        f(x_n+T)-f(x_n)
        =f(y_n)-f(y_n+S)\to0.
    \]
    因此结论对任意 $T\in\mathbb{R}$ 成立.
\end{solution}

%例题3.3.14
\begin{example}
    设 $f$ 在 $\lbrack 0,n\rbrack$ 上连续,$f(0)=f(n)$,试证:至少存在 $n$ 组不同的解 $(x,y)$,使得
    \[
        f(x)=f(y),\qquad y-x\in\mathbb{N}.
    \]
\end{example}
\exampleblank

\begin{solution}
    对 $n$ 用数学归纳法.

    当 $n=1$ 时,$f(0)=f(1)$,故 $(0,1)$ 即为一组所求解.

    假设结论对 $n=k$ 成立,现考虑 $n=k+1$.定义
    \[
        g(x)=f(x+1)-f(x),\qquad 0\leq x\leq k.
    \]
    在整数点处有
    \[
        \sum\limits_{i=0}^{k}g(i)
        =\sum\limits_{i=0}^{k}[f(i+1)-f(i)]
        =f(k+1)-f(0)=0.
    \]
    若某个 $g(i)=0$,则取 $\xi=i$;否则这些数中必有正值和负值.由于 $g$ 连续,由介值定理也可取到某个 $\xi\in\lbrack0,k\rbrack$,使
    \[
        g(\xi)=0,
    \]
    即
    \[
        f(\xi+1)=f(\xi).
    \]
    因此 $(\xi,\xi+1)$ 已是一组所求解.

    下面把区间中长度为 $1$ 的这一段``接掉''.定义 $\varphi:\lbrack0,k\rbrack\to\mathbb{R}$ 为
    \[
        \varphi(x)=
        \begin{cases}
            f(x),   & 0\leq x\leq\xi,\\
            f(x+1), & \xi<x\leq k.
        \end{cases}
    \]
    由于 $f(\xi)=f(\xi+1)$,$\varphi$ 在 $x=\xi$ 处连续,故 $\varphi\in C(\lbrack0,k\rbrack)$.又
    \[
        \varphi(0)=f(0),\qquad
        \varphi(k)=f(k+1)=f(0),
    \]
    所以由归纳假设,$\varphi$ 至少有 $k$ 组不同的解 $(u_j,v_j)$($j=1,2,\ldots,k$),满足
    \[
        \varphi(u_j)=\varphi(v_j),
        \qquad
        v_j-u_j\in\mathbb{N}.
    \]

    定义
    \[
        J(t)=
        \begin{cases}
            t,   & t\leq\xi,\\
            t+1, & t>\xi.
        \end{cases}
    \]
    则 $\varphi(t)=f(J(t))$,且 $J$ 严格递增.于是每一组 $(u_j,v_j)$ 都对应于原函数的一组解
    \[
        (J(u_j),J(v_j)).
    \]
    若 $u_j,v_j$ 位于 $\xi$ 的同侧,则
    \[
        J(v_j)-J(u_j)=v_j-u_j\in\mathbb{N};
    \]
    若 $u_j\leq\xi<v_j$,则
    \[
        J(v_j)-J(u_j)=v_j-u_j+1\in\mathbb{N}.
    \]
    因为 $J$ 为单射,这 $k$ 组解仍彼此不同.同时它们均不同于 $(\xi,\xi+1)$:若某组跨过 $\xi$,其右端经 $J$ 后严格大于 $\xi+1$;若不跨过 $\xi$,则不可能同时得到这两个端点.

    因而原函数 $f$ 至少有 $k+1$ 组不同的所求解.归纳完成.
\end{solution}

\subsection{一致连续性}

用定义证明 $f$ 在 $I$ 上一致连续,通常的方法是设法证明 $f$ 在 $I$ 满足 Lipschitz 条件:
\[
    |f(x')-f(x'')|\leq L|x'-x''|, \qquad \forall x',x''\in I,
\]
其中 $L>0$ 为某一常数.特别地,若 $f$ 在 $I$ 上存在有界导函数,则 $f$ 在 $I$ 满足 Lipschitz 条件.

而 $f(x)$ 在 $I$ 上非一致连续 $\Longleftrightarrow$ 存在 $\varepsilon_0>0$,对任意 $\dfrac{1}{n}>0$,存在 $x'_n,x''_n\in I$,虽然 $|x'_n-x''_n|<\dfrac{1}{n}$,
但是 $|f(x'_n)-f(x''_n)|\geq\varepsilon_0$.
因此,常见的证非一致连续方法是取出两个充分靠近的数列($n\to\infty$),使得它们对应函数值并不充分接近.

%例题3.3.15
\begin{example}
    判断下列函数在指定区间上的一致连续性:
    \begin{enumerate}
        \item $f(x)=\sin^2x$,$\lbrack 0,+\infty)$;
        \item $f(x)=\sin x^2$,$\lbrack 0,+\infty)$;
        \item $f(x)=\sqrt{x}$,$\lbrack 0,+\infty)$;
        \item $f(x)=\dfrac{1}{x}$,$(0,1)$;
        \item $g(x)=x^2$,$(1,+\infty)$.
    \end{enumerate}
\end{example}
\exampleblank

\begin{solution}
    \leavevmode
    \begin{enumerate}
        \item 有
              \[
                  f'(x)=\sin2x,\qquad |f'(x)|\leq1.
              \]
              由 Lagrange 中值定理,
              \[
                  |f(x)-f(y)|\leq|x-y|,
              \]
              故 $f$ 在 $\lbrack0,+\infty)$ 上满足 Lipschitz 条件,从而一致连续.

        \item 取
              \[
                  x_n=\sqrt{2n\pi},
                  \qquad
                  y_n=\sqrt{2n\pi+\dfrac{\pi}{2}}.
              \]
              则
              \[
                  |y_n-x_n|
                  =\dfrac{\pi}
                  {2\left(\sqrt{2n\pi+\dfrac{\pi}{2}}+\sqrt{2n\pi}\right)}
                  \longrightarrow0,
              \]
              但
              \[
                  f(x_n)=0,\qquad f(y_n)=1.
              \]
              因而 $|f(y_n)-f(x_n)|=1$,故 $f$ 在 $\lbrack0,+\infty)$ 上非一致连续.

        \item 不妨设 $x\geq y\geq0$.则
              \[
                  (\sqrt{x}-\sqrt{y})^2
                  =x+y-2\sqrt{xy}
                  \leq x-y,
              \]
              从而
              \[
                  |\sqrt{x}-\sqrt{y}|
                  \leq\sqrt{|x-y|}.
              \]
              给定 $\varepsilon>0$,取 $\delta=\varepsilon^2$,当 $|x-y|<\delta$ 时有 $|\sqrt{x}-\sqrt{y}|<\varepsilon$.故 $f$ 一致连续.

        \item 取
              \[
                  x_n=\dfrac{1}{n},
                  \qquad
                  y_n=\dfrac{1}{n+1}.
              \]
              则 $|x_n-y_n|\to0$,而
              \[
                  \left|\dfrac{1}{x_n}-\dfrac{1}{y_n}\right|=1.
              \]
              故 $f(x)=\dfrac{1}{x}$ 在 $(0,1)$ 上非一致连续.

        \item 取
              \[
                  x_n=n,
                  \qquad
                  y_n=n+\dfrac{1}{n}.
              \]
              则 $|x_n-y_n|=\dfrac{1}{n}\to0$,而
              \[
                  |y_n^2-x_n^2|
                  =2+\dfrac{1}{n^2}\longrightarrow2.
              \]
              故 $g(x)=x^2$ 在 $(1,+\infty)$ 上非一致连续.
    \end{enumerate}
\end{solution}

\begin{theorem}[\markstar]
    \leavevmode
    \begin{enumerate}
        \item 设 $f(x)$ 在 $(a,b)$ 上连续,则
              \[
                  f(x)\text{ 在 }(a,b)\text{ 上一致连续}
                  \Longleftrightarrow
                  \lim\limits_{x\to a^+}f(x),\ \lim\limits_{x\to b^-}f(x)\text{ 存在};
              \]
        \item 设 $f\in C(\lbrack a,+\infty))$,$\displaystyle\lim\limits_{x\to+\infty}f(x)=A$(有限),则 $f$ 在 $\lbrack a,+\infty)$ 上一致连续;
        \item 设 $f\in C(( -\infty,+\infty))$,若存在极限
              \[
                  \displaystyle\lim\limits_{x\to-\infty}f(x)=l, \qquad \lim\limits_{x\to+\infty}f(x)=L,
              \]
              则 $f(x)$ 在 $(-\infty,+\infty)$ 上一致连续.
    \end{enumerate}
\end{theorem}

\begin{remark}

    (1) 与 (2)(3) 的区别在于,无穷区间下必要性不再成立,可考虑反例 $f(x)=x,g(x)=\sin x$ 在 $(-\infty,+\infty)$ 上一致连续,但在端点 $\pm\infty$ 无极限.
\end{remark}

%例题3.3.16
\begin{example}
    判断下列函数在指定区间上的一致连续性:
    \begin{enumerate}
        \item $f(x)=\ln x$,$(0,1\rbrack$;
        \item $f(x)=\dfrac{\sin x}{x}$,$(0,1\rbrack$;
        \item $f(x)=e^x\cos\dfrac{1}{x}$,$(0,1\rbrack$;
        \item $f(x)=\cos x\cdot\cos\dfrac{\pi}{x}$,$(0,1)$;
        \item $f(x)=x\sin\dfrac{1}{x}$,$(0,+\infty)$;
        \item $f(x)=\dfrac{x+2}{x+1}\sin\dfrac{1}{x}$,$(0,1)$ 及 $(1,+\infty)$.
    \end{enumerate}
\end{example}
\exampleblank

\begin{solution}
    \leavevmode
    \begin{enumerate}
        \item 有
              \[
                  \lim\limits_{x\to0^+}\ln x=-\infty.
              \]
              若 $f$ 在 $(0,1\rbrack$ 上一致连续,则其限制在 $(0,1)$ 上也一致连续.由定理 3.3.1(1),$\displaystyle\lim\limits_{x\to0^+}f(x)$ 应为有限极限,矛盾.故 $f$ 非一致连续.

        \item 由于
              \[
                  \lim\limits_{x\to0^+}\dfrac{\sin x}{x}=1,
              \]
              可定义 $f(0)=1$,使 $f$ 连续延拓到 $\lbrack0,1\rbrack$.由 Cantor 定理,延拓后的函数在 $\lbrack0,1\rbrack$ 上一致连续,故原函数在 $(0,1\rbrack$ 上一致连续.

        \item 当 $x\to0^+$ 时,$e^x\to1$,而 $\cos\dfrac{1}{x}$ 无极限,故
              \[
                  \lim\limits_{x\to0^+}e^x\cos\dfrac{1}{x}
              \]
              不存在.若原函数一致连续,其在 $(0,1)$ 上的限制由定理 3.3.1(1) 应在 $0^+$ 处有有限极限,矛盾.故非一致连续.

        \item 取
              \[
                  x_n=\dfrac{1}{2n},
                  \qquad
                  y_n=\dfrac{1}{2n+1}.
              \]
              则 $x_n,y_n\to0^+$,且
              \[
                  \cos\dfrac{\pi}{x_n}=1,
                  \qquad
                  \cos\dfrac{\pi}{y_n}=-1.
              \]
              又 $\cos x_n,\cos y_n\to1$,所以函数在 $0^+$ 处无极限.由定理 3.3.1(1),它在 $(0,1)$ 上非一致连续.

        \item 当 $x\to0^+$ 时,
              \[
                  \left|x\sin\dfrac{1}{x}\right|\leq x\longrightarrow0,
              \]
              故可在 $0$ 处连续补值为 $0$,从而函数在 $(0,1\rbrack$ 上一致连续.另一方面,
              \[
                  \lim\limits_{x\to+\infty}x\sin\dfrac{1}{x}=1.
              \]
              由定理 3.3.1(2),函数在 $\lbrack1,+\infty)$ 上一致连续.综上可知函数在整个 $(0,+\infty)$ 上一致连续.

        \item 在 $(0,1)$ 上,因 $\dfrac{x+2}{x+1}\to2$ 而 $\sin\dfrac{1}{x}$ 在 $0^+$ 处振荡,故函数在 $0^+$ 处无极限.由定理 3.3.1(1),它在 $(0,1)$ 上非一致连续.

              对 $(1,+\infty)$,把同一表达式视为定义在 $\lbrack1,+\infty)$ 上的连续函数.由于
              \[
                  \lim\limits_{x\to+\infty}
                  \dfrac{x+2}{x+1}\sin\dfrac{1}{x}=0,
              \]
              由定理 3.3.1(2),它在 $\lbrack1,+\infty)$ 上一致连续,从而在 $(1,+\infty)$ 上一致连续.
    \end{enumerate}
\end{solution}

\begin{proposition}[\markstar]
    设 $f(x)$ 在 $\lbrack a,+\infty)$ 上一致连续,且 $g\in C(\lbrack a,+\infty))$,若有
    \[
        \displaystyle\lim\limits_{x\to+\infty}[f(x)-g(x)]=0,
    \]
    则 $g(x)$ 在 $\lbrack a,+\infty)$ 上一致连续.
\end{proposition}

\begin{corollary}[\markstar]
    设 $f\in C(\lbrack a,+\infty))$,且 $f(x)$ 有渐近线 $y=ax+b$,则 $f(x)$ 在 $\lbrack a,+\infty)$ 上一致连续.
\end{corollary}

\begin{remark}
    此即命题 3.3.1 中取 $f(x)=ax+b$.
\end{remark}

%例题3.3.17
\begin{example}
    判断下列函数 $f(x)$ 在指定区间上的一致连续性:
    \begin{enumerate}
        \item $f(x)=\dfrac{x^3}{x+1}\sin\dfrac{1}{x}$,$(0,+\infty)$;
        \item $f(x)=\sqrt{\dfrac{x^4+1}{x^2+1}}$,$(-\infty,+\infty)$;
        \item $f(x)=\sqrt[3]{x^3-x^2-x+1}$,$(-\infty,+\infty)$.
    \end{enumerate}
\end{example}
\exampleblank

\begin{solution}
    \leavevmode
    \begin{enumerate}
        \item 当 $x\to0^+$ 时,
              \[
                  \left|\dfrac{x^3}{x+1}\sin\dfrac{1}{x}\right|
                  \leq\dfrac{x^3}{x+1}\longrightarrow0,
              \]
              故可在 $0$ 处连续补值为 $0$,于是函数在 $(0,1\rbrack$ 上一致连续.

              当 $x\to+\infty$ 时,由 $\displaystyle\sin\dfrac{1}{x}=\dfrac{1}{x}-\dfrac{1}{6x^3}+O\left(\dfrac{1}{x^5}\right)$,
              \[
                  \dfrac{x^3}{x+1}\sin\dfrac{1}{x}
                  =x-1+O\left(\dfrac{1}{x}\right).
              \]
              因而直线 $y=x-1$ 是其在 $+\infty$ 处的渐近线.由命题 3.3.1 的推论,函数在 $\lbrack1,+\infty)$ 上一致连续.两部分在 $x=1$ 处拼接,故原函数在 $(0,+\infty)$ 上一致连续.

        \item 记
              \[
                  F(x)=\sqrt{\dfrac{x^4+1}{x^2+1}}.
              \]
              函数 $F$ 为偶函数.当 $x\to+\infty$ 时,
              \[
                  F(x)-x
                  =\dfrac{\dfrac{x^4+1}{x^2+1}-x^2}{F(x)+x}
                  =\dfrac{1-x^2}{(x^2+1)(F(x)+x)}
                  \longrightarrow0.
              \]
              函数 $x$ 在 $\lbrack0,+\infty)$ 上一致连续,故由命题 3.3.1,$F$ 在 $\lbrack0,+\infty)$ 上一致连续.利用 $F$ 的偶性可知它在 $(-\infty,0\rbrack$ 上也一致连续.两侧在 $0$ 处拼接,因此 $F$ 在 $\mathbb{R}$ 上一致连续.

        \item 记
              \[
                  F(x)=\sqrt[3]{x^3-x^2-x+1}.
              \]
              当 $|x|\to\infty$ 时,
              \[
                  \begin{aligned}
                      F(x)
                      &=x\left(1-\dfrac{1}{x}-\dfrac{1}{x^2}+\dfrac{1}{x^3}\right)^{\frac{1}{3}}\\
                      &=x-\dfrac{1}{3}+O\left(\dfrac{1}{x}\right).
                  \end{aligned}
              \]
              因而在 $+\infty$ 与 $-\infty$ 两端,$F(x)$ 都渐近于一次函数
              \[
                  y=x-\dfrac{1}{3}.
              \]
              取充分大的 $A>0$.由命题 3.3.1 的推论,$F$ 在 $\lbrack A,+\infty)$ 上一致连续;对函数 $G(t)=F(-t)$ 应用同一推论,可知 $F$ 在 $(-\infty,-A\rbrack$ 上一致连续.而 $F$ 在紧区间 $\lbrack-A,A\rbrack$ 上连续,由 Cantor 定理一致连续.将这三个相邻区间依次拼接,得到 $F$ 在 $\mathbb{R}$ 上一致连续.
    \end{enumerate}
\end{solution}

\begin{proposition}[\markstar\ 阶的估计]
    设 $f(x)$ 在 $\mathbb{R}$ 上一致连续,则存在 $a,b>0$,使得
    \[
        |f(x)|\leq a|x|+b.
    \]
\end{proposition}

\begin{remark}
    对比命题 3.3.1 及其推论与命题 3.3.2.前者说明:若连续函数 $f$ 在无穷远处与某个一次函数 $ax+b$ 之差趋于 $0$,则 $f$ 一致连续;后者说明:一致连续函数的增长至多为线性增长,即存在 $a,b>0$,使 $|f(x)|\leq a|x|+b$.
\end{remark}

\begin{proposition}[\markstar\ 四则运算]
    设 $f(x),g(x)$ 在区间 $I$ 上一致连续,则:
    \begin{enumerate}
        \item 对任意 $k\in\mathbb{R}$,$kf(x)$ 与 $f(x)\pm g(x)$ 在 $I$ 上一致连续;
        \item $I$ 为有限区间时,$f(x)g(x)$ 在 $I$ 上一致连续;$I$ 为无穷区间时,$f(x)g(x)$ 在 $I$ 上不一定一致连续;
        \item $f(x)\neq0$ 时,$\dfrac{1}{f(x)}$ 与 $\dfrac{g(x)}{f(x)}$ 在 $I$ 上不一定一致连续.
    \end{enumerate}
\end{proposition}

\begin{proposition}[\markstar\ 函数复合]
    设 $f(x)$ 在 $I_1$ 上一致连续,$g(x)$ 在 $I_2$ 上一致连续,且 $f(x)$ 的值域包含在 $I_2$ 中,那么 $g(f(x))$ 在 $I_1$ 上也一致连续.
\end{proposition}

%例题3.3.18
\begin{example}
    设 $f$ 在 $\lbrack 0,+\infty)$ 上满足 Lipschitz 条件,证明:$f(x^\alpha)\,(0<\alpha<1)$ 在 $\lbrack 0,+\infty)$ 上一致连续.
\end{example}
\exampleblank

\begin{solution}
    设 $f$ 的 Lipschitz 常数为 $L$.对 $x,y\geq0$ 且 $0<\alpha<1$,有 $|x^\alpha-y^\alpha|\leq|x-y|^\alpha$.
    因而
    \[
        |f(x^\alpha)-f(y^\alpha)|\leq L|x^\alpha-y^\alpha| \leq L|x-y|^\alpha.
    \]
    右边只依赖于 $|x-y|$,并在其趋于零时趋于零,故复合函数一致连续.
\end{solution}

%例题3.3.19
\begin{example}[区间合并]
    \begin{enumerate}
        \item 若函数 $f$ 在区间 $(a,b\rbrack$ 和 $\lbrack b,c)$ 上分别一致连续,证明:$f$ 在 $(a,c)$ 上一致连续;
        \item 若函数 $f$ 在区间 $(a,b)$ 和 $\lbrack b,c)$ 上分别一致连续,此时 $f$ 在 $(a,c)$ 上是否一致连续?
        \item 证明:函数 $f(x)=\dfrac{|\sin x|}{x}$ 在区间 $(-1,0)$ 和 $(0,1)$ 均一致连续,但在 $(-1,0)\cup(0,1)$ 上非一致连续.
    \end{enumerate}
\end{example}
\exampleblank

\begin{solution}
    第 (1) 问,对给定 $\varepsilon$ 分别取两侧的一致连续模.若 $x,y$ 位于同侧则直接使用;若 $x<b<y$,则 $|f(x)-f(y)|\leq|f(x)-f(b)|+|f(b)-f(y)|$,
    故在 $(a,c)$ 上一致连续.

    第 (2) 问不一定.取 $f=0$ 于 $(a,b)$,$f=1$ 于 $\lbrack b,c)$,两段分别一致连续,但合并后不连续.

    第 (3) 问,该函数在两侧分别可连续延拓到闭区间,故分别一致连续.但取 $x_n=\dfrac{1}{n},y_n=-\dfrac{1}{n}$,则 $|x_n-y_n|\to0$,而
    \[
        \left|\dfrac{|\sin x_n|}{x_n}-\dfrac{|\sin y_n|}{y_n}\right|\to2,
    \]
    故并集上非一致连续.
\end{solution}

%例题3.3.20
\begin{example}
    设 $f(x)$ 在 $\lbrack 0,+\infty)$ 上一致连续,且有
    \[
        \lim\limits_{n\to\infty}f(x+n)=0,\qquad \forall x\in\lbrack 0,+\infty).
    \]
    则
    \[
        \lim\limits_{x\to+\infty}f(x)=0.
    \]
\end{example}
\exampleblank

\begin{solution}
    给定 $\varepsilon>0$.由 $f$ 在 $\lbrack0,+\infty)$ 上一致连续,存在 $\delta>0$,使对任意 $u,v\geq0$,当 $|u-v|<\delta$ 时,
    \[
        |f(u)-f(v)|<\dfrac{\varepsilon}{2}.
    \]
    取正整数 $m$,使 $\dfrac{1}{m}<\delta$,并把 $\lbrack0,1\rbrack$ 分成有限个分点
    \[
        t_i=\dfrac{i}{m},\qquad i=0,1,\ldots,m.
    \]
    对每个固定的 $t_i$,由题设
    \[
        f(t_i+n)\longrightarrow0\qquad(n\to\infty).
    \]
    由于分点只有有限个,存在同一个正整数 $N$,使对所有 $n\geq N$ 及所有 $i=0,1,\ldots,m$,
    \[
        |f(t_i+n)|<\dfrac{\varepsilon}{2}.
    \]

    任取 $x\geq N+1$,写成
    \[
        x=n+t,\qquad n=\lfloor x\rfloor,\qquad 0\leq t<1.
    \]
    此时 $n\geq N$.可取某个分点 $t_i$,使
    \[
        |t-t_i|\leq\dfrac{1}{m}<\delta.
    \]
    因而
    \[
        |x-(n+t_i)|=|t-t_i|<\delta,
    \]
    由一致连续性,
    \[
        |f(x)-f(n+t_i)|<\dfrac{\varepsilon}{2}.
    \]
    再结合 $|f(n+t_i)|<\dfrac{\varepsilon}{2}$,得到
    \[
        |f(x)|
        \leq|f(x)-f(n+t_i)|+|f(n+t_i)|
        <\varepsilon.
    \]
    故
    \[
        \lim\limits_{x\to+\infty}f(x)=0.
    \]
\end{solution}

%例题3.3.21
\begin{example}
    设 $f(x)$ 在区间 $I$ 上有定义,记
    \[
        w_f(\delta) =\sup\limits_{\substack{x',x''\in I\\|x'-x''|<\delta}} |f(x')-f(x'')|,
    \]
    则 $f$ 在 $I$ 上一致连续的充要条件是
    \[
        \displaystyle\lim\limits_{\delta\to0^+}w_f(\delta)=0.
    \]
\end{example}
\exampleblank

\begin{solution}
    若 $f$ 一致连续,则对任意 $\varepsilon>0$ 存在 $\delta_0>0$,使 $|x'-x''|<\delta_0$ 时 $|f(x')-f(x'')|<\varepsilon$,故 $w_f(\delta_0)\leq\varepsilon$.又 $w_f$ 随 $\delta$ 单调不减,故 $\delta\to0^+$ 时 $w_f(\delta)\to0$.

    反之,若 $w_f(\delta)\to0$,给定 $\varepsilon>0$,取 $\delta>0$ 使 $w_f(\delta)<\varepsilon$.当 $|x'-x''|<\delta$ 时,根据上确界的定义有 $|f(x')-f(x'')|\leq w_f(\delta)<\varepsilon$,
    故 $f$ 一致连续.
\end{solution}
